\section{Introduction and the exact search relation}\label{sec:introduction}

The class PPAD organizes total search problems whose existence proofs are
based on the parity of endpoints in a directed graph~\cite{Papadimitriou1994}. Its canonical problem,
End-of-Line, asks for an endpoint other than a specified source, when the
graph is described by predecessor and successor circuits. Generalized
circuits provide another useful language for this class: a solution assigns
real values in $[0,1]$ to wires and approximately satisfies local arithmetic,
comparison, and Boolean constraints. The generalized-circuit intermediary
was developed by Chen, Deng, and Teng~\cite[Section~4.2]{CDT2007Preprint};
arithmetic, comparison, and Boolean gadgets also underlie the equilibrium
reductions of Daskalakis, Goldberg, and
Papadimitriou~\cite[Section~5.1]{DGP2009}. Rubinstein established
constant-error hardness with every gate required to be
satisfied~\cite[Theorem~2]{Rubinstein2018}. These results motivate the
circuit language used here, with the precise normalization and thresholds
specified below.

Babichenko, Papadimitriou, and Rubinstein formulated the PCP-for-PPAD
conjecture as a quasilinear reduction from End-of-Line to generalized
circuits in which both the allowed numerical error and the allowed fraction
of violated gates are positive constants~\cite[Conjecture~2]{BPR2016}.
The failure fraction is the essential robustness requirement: the decoder
must succeed even if an adversary chooses which constraints to violate.
The quasilinear size requirement is a separate restriction on the reduction.
Ordinary polynomial-time PPAD-hardness, even with constant numerical error,
does not provide it. For example, hardness for every numerical error below $1/10$ is known
for the standard generalized-circuit problem with every gate required to
be satisfied~\cite[Theorem~4.1]{DFHM2024}. Subsequent uses of PCP-for-PPAD also
consider the weaker conjecture asserting ordinary PPAD-hardness with a
positive failure fraction; the distinction from the original quasilinear
formulation is explicit in~\cite[Section~2.2.2]{Fisher2026}.

Rubinstein's 2016 work combined encoded-path geometry, a locally
evolving version of End-of-Line, and a holographic proof system whose
certificates can be constructed and decoded through restricted
access~\cite[Sections~3, 5, and~6]{Rubinstein2016Settling}. It established
quasipolynomial hardness for constant-error two-player Nash under the
exponential-time hypothesis for PPAD, while leaving the PCP-for-PPAD
conjecture open and identifying composition as a remaining
question~\cite[Theorem~1.2 and Section~1.3]{Rubinstein2016Settling}.
More recently, Golowich derives sharper two-player Nash lower bounds
for vanishing approximation error under PCP for PPAD together with
the exponential-time hypothesis for PPAD~\cite[Theorem~1.1]{Golowich2026}.
These time lower bounds retain their additional exponential-time
hypothesis; the result here concerns the conjectured reduction itself.

We give a positive resolution of the quasilinear PCP-for-PPAD conjecture.
Our size bound counts the complete finite output, including connection
indices, every auxiliary gate, and the bits of all rational parameters.
The gate conventions below are the standard bounded-rational conventions
of~\cite{DFHM2024,Fisher2026}; the totality and witness-size statement is
proved explicitly. When restricted to the nearly linear predecessor and
successor circuits in the original formulation, the theorem gives the
conjectured quasilinear reduction there as well.\footnote{The original
Definition~1 of~\cite{BPR2016} allows $x=0^n$ in its first solution
disjunct $P(S(x))\ne x$. If $P(S(0^n))\ne0^n$, this is an immediately
computable solution. Otherwise the instance satisfies the additional
source promise used below. Thus this normalization requires only circuit
evaluation and a trivial branch of the reduction: output a fixed one-gate
Constant circuit and have the decoder return $0^n$.}

\subsection{The source and target problems}

We fix conventional explicit binary encodings of circuit gate lists,
connection indices, input and output lists, and rational parameters.
Integers are encoded in binary with delimiters. All size bounds concern
these complete encodings. Changing between the usual such explicit
encodings changes only constant or logarithmic factors.

\begin{definition}[End-of-Line]\label{def:eol}
An instance $I$ consists of two explicitly listed, acyclic, bounded-fanin
Boolean circuits
\[
 S,P:\{0,1\}^n\longrightarrow\{0,1\}^n,\qquad n\geq1,
\]
using AND, OR, and NOT gates, with unrestricted fan-out, and satisfying
\[
 P(0^n)=0^n,\qquad S(0^n)\ne0^n,\qquad P(S(0^n))=0^n.
\]
A solution is a string $v\ne0^n$ such that
\begin{equation}\label{eq:end-of-line-solution}
 P(S(v))\ne v\quad\text{or}\quad S(P(v))\ne v.
\end{equation}
The input length $N\geq2$ includes the complete descriptions of both
circuits and their input and output lists.
\end{definition}

\begin{definition}[Generalized circuit]\label{def:gcircuit}
A generalized circuit $G=(V,T)$ has an explicitly listed finite set of
nodes $V$ and a nonempty set of gates $T$. Each gate has one output and at
most two inputs. A node is the output of at most one gate. Directed cycles
and unrestricted fan-out are permitted. Every gate parameter $\alpha$ is
a binary-encoded rational in $[0,1]$.

For an assignment $x:V\to[0,1]$, write $x_a,x_b$ for input values and
$x_z$ for the output. The five arithmetic gate types are
$\eps$-satisfied when the corresponding condition holds:
\begin{align*}
 \mathrm{Constant}(\alpha;z)&:&\quad \abs{x_z-\alpha}&\leq\eps,\\
 \mathrm{Scale}(\alpha;a;z)&:& \abs{x_z-\alpha x_a}&\leq\eps,\\
 \mathrm{Copy}(a;z)&:& \abs{x_z-x_a}&\leq\eps,\\
 \mathrm{Add}(a,b;z)&:& \abs{x_z-\min\{x_a+x_b,1\}}&\leq\eps,\\
 \mathrm{Subtract}(a,b;z)&:& \abs{x_z-\max\{x_a-x_b,0\}}&\leq\eps.
\end{align*}
The comparison gate $\mathrm{Less}(a,b;z)$ has the two implications
\[
 x_a<x_b-\eps\ \Longrightarrow\ x_z\geq1-\eps,
 \qquad
 x_a>x_b+\eps\ \Longrightarrow\ x_z\leq\eps.
\]
The Boolean gates have the following implications:
\begin{align*}
 \mathrm{Or}(a,b;z):\qquad
 (x_a\geq1-\eps\ \text{or}\ x_b\geq1-\eps)&\Longrightarrow x_z\geq1-\eps,\\
 (x_a\leq\eps\ \text{and}\ x_b\leq\eps)&\Longrightarrow x_z\leq\eps;\\[2pt]
 \mathrm{And}(a,b;z):\qquad
 (x_a\geq1-\eps\ \text{and}\ x_b\geq1-\eps)&\Longrightarrow x_z\geq1-\eps,\\
 (x_a\leq\eps\ \text{or}\ x_b\leq\eps)&\Longrightarrow x_z\leq\eps;\\[2pt]
 \mathrm{Not}(a;z):\qquad
 x_a\leq\eps&\Longrightarrow x_z\geq1-\eps,\\
 x_a\geq1-\eps&\Longrightarrow x_z\leq\eps.
\end{align*}
For these four gate types, every applicable implication must hold. If no
antecedent holds, any output in $[0,1]$ is permitted. In particular, the
comparison thresholds are strict and the Boolean thresholds are weak.

Let $M$ be the complete binary encoding length of $G$. For fixed
$\eps,\delta>0$ and a fixed polynomial $p$, an acceptable assignment is
rational, has total binary encoding length at most $p(M)$, and
$\eps$-satisfies at least $(1-\delta)|T|$ gates. Each gate has equal weight;
at most $\lfloor\delta|T|\rfloor$ gates may fail. Failed gates may violate
their relations arbitrarily, but all their nodes still have one globally
consistent value in $[0,1]$.
\end{definition}

The Boolean predicates make the error parameter nonmonotone: increasing
$\eps$ can activate an antecedent that was previously vacuous. This issue
is analyzed by Schuldenzucker and Seuken~\cite[Proposition~1]{SchuldenzuckerSeuken2019}.
Our gate arguments use one actual tolerance throughout. Later upper bounds
on that tolerance ensure explicit margins; they do not reclassify a
satisfied Boolean gate at a larger tolerance.

\subsection{Main theorem}

\begin{samepage}
\begin{theorem}[Quasilinear PCP for PPAD]\label{thm:main}
There are fixed rational constants
\[
 0<\eps<\tfrac1{10},\qquad 0<\delta<1,
\]
a fixed positive integer $A$, a fixed nonnegative integer $b$, a fixed
polynomial $p$ with nonnegative integer coefficients, and deterministic
algorithms $R,D$ with the following properties.
\begin{enumerate}[label=\textup{(\roman*)}]
\item Every generalized circuit in Definition~\ref{def:gcircuit} has an
acceptable rational assignment of encoding length at most $p(M)$.
\item For every End-of-Line instance $I$ of length $N$, the output $R(I)$
is a generalized circuit whose complete encoding length satisfies
\[
 M\leq A N\bigl\lceil\log_2(N+2)\bigr\rceil^b.
\]
\item The construction $R$ runs in time polynomial in $N$, with one fixed
polynomial bound.
\item For every acceptable assignment $x$ of $R(I)$, the algorithm $D(I,x)$
returns a solution of $I$ in time polynomial in $N$ plus the encoding
length of $x$, with one fixed polynomial bound.
\end{enumerate}
All constants, polynomials, and algorithms are chosen once. They use no
advice depending on the input length and no randomized construction.
\end{theorem}
\end{samepage}

The conclusion about $D$ is universal over the locations and values of
the failed gates. In fact, the soundness argument applies to arbitrary
real assignments meeting the gate-failure budget; rationality is needed
for the finite search relation and its computational decoder. The totality
proof supplies an assignment satisfying every gate on a fixed rational
mesh. Neither tolerance tends to zero with $N$.

\subsection{Proof strategy: encoded paths and robust local computation}

We begin with the geometric target of the reduction. Replace the
End-of-Line graph by a locally evolving directed graph with a known
source, no isolated vertices, and the same endpoint-decoding property.
Represent each vertex $w$ by a separated binary word $E_w\in\{0,1\}^m$.
In the box $[-1,2]^{4m}$, viewed as four blocks of $m$ coordinates,
each edge $u\to v$ traces
a polygon that first changes a copy of $E_u$ to $E_v$, changes a repeated
control bit, updates the other copy, and resets the control bit. Rounded
corners join these polygons into paths and cycles. An additional segment
enters the known source, and a separate field prescription handles its
far end. Every other path endpoint encodes an End-of-Line answer.

Section~\ref{sec:geometry} constructs a bounded displacement field $v$
on this box. Close to a path it uses the directed tangent and the inward
normal; farther away it passes to a fixed direction. For a sufficiently
small fixed step $\beta>0$, the clipped move
$\clip(X+\beta v(X))-X$ has norm at least a fixed positive multiple of
$\beta$ outside a fixed neighborhood of the desired endpoints. Here
clipping keeps each coordinate in the box, and the norm is normalized
by the word length. Thus an accurate enough approximate fixed point
identifies an endpoint. The proof must also provide local numerical
access to the field: it cannot inspect the entire encoded graph when
computing one coordinate of $v$.

This use of encoded paths and restricted access has a close predecessor
in Rubinstein's holographic proof system~\cite[Sections~3, 5, and~6]{Rubinstein2016Settling}.
The present reduction needs two further properties simultaneously:
quasilinear total storage and robustness to an arbitrary fixed fraction
of failed gates. The two scales of the construction address these
requirements as follows.

\paragraph{A quasilinear outer word.}
Sections~\ref{sec:macrograph}--\ref{sec:outer-representation} construct
canonical representations
\[
 R_w=(s_k)_{k\in\Gamma},\qquad
 s_k\in\{0,1\}^L,\qquad
 |\Gamma|=N(\log N)^{O(1)},\qquad L=(\log N)^{O(1)}.
\]
A bounded number of symbol positions suffices to test local validity,
propose a neighboring symbol, or perform one correction step. Each such
operation has Boolean complexity polynomial in $L$, with an exponent
fixed before the later correction accuracy. Distinct valid words have
constant symbol distance, and a word within a fixed symbol-error radius
can be decoded globally in polynomial time.

Begin with the sparse transcript graph of Section~\ref{sec:macrograph},
whose vertices record partial computations of the predecessor and successor
circuits and whose edges change one tape position and its occupation bit.
Transcript validity is expressed by bounded-degree polynomial equations in
the tape entries. Sampling an individual equation could miss one false
equation. Sections~\ref{sec:outer-codes} and~\ref{sec:outer-programs}
instead encode polynomial remainders recording the equation residuals,
and retain certificates for the intermediate products and linear
computations that produce them. On exact codewords, a nonzero remainder
differs from zero at a constant fraction of positions. The code's commuting
row and column constraints also permit whole-row or whole-column comparisons
that give a constant gap for an incorrect certificate identity. All stored
components are tested together on one fixed list of neighboring positions,
so more certificates enlarge a symbol without diluting the gap.

One tape edit can change all these encoded arrays. Before making an allowed
edit, Section~\ref{sec:outer-representation} prepares each retained array's
finite difference: its value on the edited tape minus its value on the old
tape. Certificates introduced by this preparation need their own increments
as well. Each new increment has lower assigned polynomial degree in the
logical tape entries, so this recursion has bounded depth. Once preparation
is complete, adding the increments updates the whole retained tuple to its
canonical value on the new tape; the preparation can then be reversed.

The remaining storage issue is the long linear computation needed to prepare
an increment. It is carried out one small group of Fourier frequencies at a
time: each such orbit is preserved by the encoder and has only
polylogarithmically many frequencies. A clock records progress and scratch
space is reused. Thus the preparation's elapsed time lengthens the
represented path rather than the data stored at a vertex. The resulting
words are the locally evolving $R_w$ above.

\paragraph{One fixed binary port.}
A polynomial-size robust construction applied directly to an $N$-bit
vertex would not give the required size. Instead we apply it only to
the $L$-bit outer symbols. The inner construction supplies one injective
encoding
\[
 \mathcal D:\{0,1\}^L\longrightarrow\{0,1\}^{S_0},
 \qquad S_0=\operatorname{poly}(L),
 \qquad E_w=(\mathcal D(s_k))_{k\in\Gamma}.
\]
Its encoding and positive correction radius are chosen before the
number of symbols inspected by a local function. From that fixed radius,
it supports every prescribed fixed average accuracy for requests with
bounded address density. More outer correction rounds may enlarge the number
of symbols inspected, and hence the number of input ports in the resulting
local computation. More inner correction rounds improve its requested
bit accuracy. These choices may enlarge constant factors, but they do
not change $\mathcal D$ or its radius. All polynomial factors in $L$
remain fixed powers of $\log N$.

Sections~\ref{sec:inner-algebra}--\ref{sec:inner-gadget} construct this
port and its robust decision gadget independently of the outer
reduction. For a fixed tuple $s=(s_1,\ldots,s_k)$ of $L$-bit inputs
to a Boolean function $f$, the inner graph is one computation path with
blank work at its source and output $f(s)$ at its sink. Polynomial charts encode its tapes and
certificates. Their dependencies are kept separate for each read-only
input, including throughout the finite-difference closure. Consequently
the input-dependent parts of the source and its first successor have
the same formulas in every later invocation. A one-input template at
these two labels defines $\mathcal D(s_i)$ for one input $s_i$; its bits then supply the
source data for an actual multi-input computation.

The binary representation in Section~\ref{sec:inner-binary} adds
point tables, line tables, and tables for the required products.
Correction keeps both addresses in each product request and preserves
the required bound on their joint distribution. Its error bound is averaged over requested coordinates;
arbitrarily accurate correction at every individual coordinate is
neither asserted nor needed. At the inner scale, the geometric
construction forces a feedback state near the unique sink, where a
positive fraction of encoded coordinates reveal $f(s)$. These become the
many designated outputs of the decision gadget.

\paragraph{From local decisions to an endpoint under gate failures.}
Section~\ref{sec:outer-composition} uses the outer tester and correction
routines to supply the field's local information. Two rounded blocks
serve as possible vertex encodings. On good validation decisions, each
accepted seed has a unique reference vertex; corrected calls supply its incident
flags and the requested bits of its own and neighboring encodings.
The geometric estimator reserves a bounded list of possible pieces and
estimates their statistics before selecting a numerical winner.
This order makes average bit accuracy sufficient even when the selected
piece depends on the output coordinate.

The generalized circuit has many physical replicas of each geometric
coordinate. Write $X_{ir}$ for replica $r$ of component $i$, and
$\bar X_i$ for their mean. A fixed-degree expander average $A$ couples
the replicas. Each output is produced by a circuit that realizes the
approximate relation
\[
 X_{ir}\approx\clip\bigl((AX)_{ir}
                       +\beta\widehat v_{ir}(\bar X)\bigr),
\]
where $\widehat v$ is the bounded local field estimator. First, boundedness
of the update and the expander gap force the replicas close to their
means, without assuming that any decision is correct. This synchronization
then bounds the number of reads far from their means. Randomized
threshold offsets control ambiguous comparisons, and the gate-failure
budget controls the circuits containing a failed gate.
The randomness describes the sampling analysis: the printed circuit assigns
replicas a deterministically enumerated list of fixed tapes of size polynomial
in $L$.

These charges apply to the complete wired computation. All hypothetical
answer branches are allocated before routing, and each physical read
occurrence is routed separately with bounded load. Expensive Boolean
decision gadgets are shared by groups of output replicas. Their many
designated outputs ensure that failures concentrated on extraction
wires still spoil only a controlled fraction of the group. On reliable
outputs the update approximates the true field; on unreliable outputs
whose gates are satisfied, its bounded size still gives an error
proportional to $\beta$, apart from the numerical gate tolerance. An estimate that keeps clipping inside the
replica contraction argument then makes the mean an approximate fixed
point. The displacement lower bound forces $\bar X$ near an endpoint.
Rounding its first block and applying the global binary and outer
decoders yields the required End-of-Line answer.

The same synchronization, error charging, and clipped-mean argument is
used inside the decision gadget. At that scale every output has its own
bounded estimator circuit; the additional sharing of decision gadgets
occurs only in the outer composition. In both uses, a fixed hypothetical
answer branch is analyzed under the original random-tape law, rather
than by conditioning on that branch being executed.

\paragraph{Organization and imported inputs.}
Section~\ref{sec:outer-composition} assembles the main theorem from
the outer representation, the port interface stated there, and the
independent geometric interface. It proves universal decoding, the
complete binary size bound, and rational-witness totality.
Sections~\ref{sec:inner-algebra}--\ref{sec:inner-gadget} prove the port
interface, and Section~\ref{sec:geometry} proves both the displacement
bound and finite numerical access. Section~\ref{sec:weak-nash} gives
the consequences described below. Appendix~\ref{sec:extensions}
collects the separate full-macrograph and unrestricted-product variants;
these are not prerequisites for the fixed-port composition.

The outer code uses ternary Goldbach, fixed-generator prime-quotient
expansion, and an effective least-prime bound
\cite{Helfgott2013,BG2008,Xylouris2011}. Its simultaneous systematic
encoders, matrix twists, and joint correction guarantees are proved
here; a general locally testable code theorem such as~\cite{DEL+2022}
does not by itself supply this combination. The inner binary construction
uses the corrected scalar line test and two group-theoretic inputs
\cite{FS1995,EJZ2010,Shalom1999}; its vector replacement, affine
expansion, and average correction are proved here. The all-modulus
Margulis--Gabber--Galil bound supplies the local and replica walks
\cite{MGGAFP}. The inner fields and groups are introduced afresh and
are distinct from those of the outer construction.

\subsection{Consequences for games and markets}

Section~\ref{sec:weak-nash} applies the theorem to the conditional
reductions for weak approximate Nash and relative bimatrix Nash
\cite{BPR2016}, course allocation~\cite{BPR2016}, markets with
approximately optimal bundles~\cite{Fisher2026}, stationary equilibria
in discounted games~\cite{DGZ2023,PZO2025,FKPP2026}, and approximate
Hylland--Zeckhauser allocation~\cite{HZ2026}. For each target we state
the precise approximation and encoding conventions before its
consequence. The cited reductions provide the transfers to these
problems. We supply the needed gate-basis and fan-out conversions, track
their fixed error and failure losses, and verify the rational witness
and decoder bounds. These consequences use fixed-gap circuit hardness;
the quasilinear complete-output bound is the additional assertion of
Theorem~\ref{thm:main}.
